\capitulo{2}{Objetivos del proyecto}

El objetivo principal de este Trabajo de Fin de Grado es el diseño y desarrollo de un sistema capaz de registrar y analizar la interacción del usuario en entornos de modelado tridimensional, utilizando Blender como plataforma de referencia.

El propósito de esta herramienta es proporcionar un mecanismo que permita capturar y estructurar información del proceso de modelado de forma automática, con el fin de facilitar su análisis posterior. De este modo, se busca obtener métricas objetivas sobre el comportamiento del usuario, el uso de las herramientas y la evolución de las escenas, contribuyendo a una mejor comprensión del proceso de modelado tridimensional \cite{jankowski2014,lavioila2017}.

Este sistema pretende capturar información relevante sobre el proceso de modelado sin interferir en el flujo de trabajo del usuario, permitiendo posteriormente analizar los datos obtenidos para estudiar cómo se construyen los modelos tridimensionales, qué operaciones se realizan con mayor frecuencia y cómo evoluciona la interacción en el entorno 3D \cite{bimanual2012}.

El proyecto se centra tanto en la recopilación de datos como en su posterior análisis, proporcionando herramientas que permitan explorar la información registrada desde diferentes perspectivas, como el análisis temporal de las acciones realizadas, el estudio de la evolución de la escena o la detección de eventos relevantes de modelado. Este enfoque se alinea con las tendencias actuales en el análisis de interacción en entornos tridimensionales y sistemas de modelado avanzados \cite{malik2023}.

\section{Objetivos específicos}

Para alcanzar este objetivo general, se han definido los siguientes objetivos específicos:

\begin{itemize}
    \item Diseñar un sistema de recopilación de datos no intrusivo que funcione de manera transparente para el usuario durante el proceso de modelado 3D.
    
    \item Implementar un addon para Blender capaz de registrar información temporal, espacial y estructural de la escena y de los objetos modelados.
    
    \item Detectar de forma automática eventos relevantes de modelado, tales como duplicados, fusiones de geometría, cambios de modo de edición y modificaciones en las coordenadas UV.
    
    \item Garantizar la persistencia de los datos recopilados entre sesiones de trabajo, integrándolos de forma segura dentro del propio archivo del proyecto.
    
    \item Incorporar un mecanismo explícito de consentimiento del usuario que asegure una recopilación de datos ética y transparente.
    
    \item Validar el correcto funcionamiento del sistema mediante pruebas experimentales controladas.
    
    \item Recopilar datos reales a partir de una versión estable del sistema utilizada en sesiones de modelado.
    
    \item Desarrollar una herramienta de análisis que permita explorar y visualizar los datos recogidos desde una perspectiva temporal, espacial y estadística.
\end{itemize}

Estos objetivos permiten estructurar el desarrollo del proyecto en diferentes fases, desde la implementación del sistema de registro de datos hasta el análisis y visualización de la información obtenida a partir de su uso en entornos reales.
